\section{tcl++}\label{tcl++}

This section describes a light weight class library for creating
Tcl/Tk applications. \cite{Ousterhout94} It consists of a header file
\verb|tcl++.h|, and a main program \verb|tclmain.cc|. All you need to
do to create a Tcl/Tk application is write any application specific
Tcl commands using the NEWCMD macro, link your code with the
\verb|tclmain| code, and then proceed to write the rest of your
application as a Tcl script. The \verb|tclmain| code uses argv[1] to
get the name of a script to execute, so assuming that the executable
image part of your application is called \verb|appl.exe|, you would
start off your application script with
\begin{verbatim}
#!appl.exe
...Your Tcl/Tk code goes here...
\end{verbatim}
then you script becomes the application, directly launchable from the
shell.

\subsection{Global Variables}

The following global symbols are defined by \verb|tcl++|:
\begin{description}
\item[interp] The default interpreter used by Tcl/Tk
\item[mainWin] The window used by Tk, of type \verb|Tk_Window|, useful
for obtaining X information such as colourmaps
\item[TclError] A longjmp label used for trapping errors --- this will
be replaced by exception handling once g++ has this ability
\item[error] To flag an error, call \verb|error(char *format,...)|.
The format string is passed to sprintf, along with all remaining
arguments. The error message is attached to Tcl's error log, and then
control branches to \verb|TclError|, which returns a \verb|TCL_ERROR|
code back to Tcl.
\end{description}

\subsection{Creating a Tcl command}
To register a procedure, use the \verb|NEWCMD| macro. It takes two
arguments, the first is the name of the Tcl command to be created, the
second is the number of arguments that the Tcl command takes.
\verb|NEWCMD| does the following:
\begin{itemize}
\item lodges  the commands with the interpreter prior to the
execution of tclmain,
\item sets up the longjmp label to trap errors,
\item checks that the declared number of arguments match the actual
number (type checking is beyond the ANSI C Preprocessor!)
\item executes the contents of the macro \verb|GLOBAL_INIT_HOOK|,
which can be set to (for example) initialize any global variables from
Tcl variables.
\item declares a void returning C++ function of the same name as the
Tcl command. The arguments are passed via the argc, argv
mechanism. All that is required is the body of the function to be put
in.
\end{itemize}

Here is an example.
\begin{verbatim}
#include "tcl++.h"

#undef GLOBAL_INIT_HOOK
#define GLOBAL_INIT_HOOK init_funny();

static int funny;

void init_funny()  /* this must be declared before NEWCMD */
{
   if ( 
     Tcl_GetInt(interp, 
        Tcl_GetVar(interp,"funny",TCL_GLOBAL_ONLY), 
        &funny) 
     !=TCL_OK) error("funny not found"); 
}


NEWCMD(silly_cmd,1)
{printf("argv[1]=%s funny=%d\n", argv[1], funny);}
\end{verbatim}
This creates a Tcl command that prints out its argument, and the value
of the variable \verb|funny|.

\subsection{tclvar}

Tcl variables can be accessed by means of the tclvar class. For
example, if the programmer declares:

\verb|tclvar hello="hello"; float floatvar;|

then the variable hello can be used just like a normal C variable in
expression such as 

\verb|floatvar=hello*3.4;|.

The class allows assignment to and from \verb|double| and \verb|char*|
variables, incrementing and decrementing the variables, compound
assignment and accessing array Tcl variables by means of the C array
syntax. The is also a function that tests for the existence of a Tcl
variable.

\pagebreak[4]
The complete class definition is given by:

\begin{verbatim}
class tclvar
{ 
 public:

/* constructors */
  tclvar();
  tclvar(char *nm, char* val=NULL);
  tclvar(tclvar&);
  ~tclvar();

/* These four statements allow tclvars to be freely mixed with arithmetic 
 expressions */
  double operator=(double x);
  char* operator=(char* x);
  operator double ();
  operator char* ();

  double operator++();
  double operator++(int);
  double operator--();
  double operator--(int); 
  double operator+=(double x);
  double operator-=(double x);
  double operator*=(double x);
  double operator/=(double x);

  tclvar operator=(tclvar x); 
/* arrays can be indexed either by integers, or by strings */
  tclvar operator[](int index);
  tclvar operator[](char* index);

  friend int exists(tclvar x);
};
\end{verbatim}

\subsection{tclcmd}

a tclcmd allows Tcl commands to be executed by a simple \verb|x <<|
{\em Tcl command} syntax. The commands can be stacked, or
accumulated. Each time a linefeed is obtained, the command is
evaluated.

\pagebreak[4]
For example:
\begin{verbatim}
tclcmd cmd;

cmd << "set a 1\n puts stdout $a";
cmd << "for {set i 0} {$i < 10} {incr i}"
cmd << "{set a [expr $a+$i]; puts stdout $a}\n";
\end{verbatim}

The behaviour of this command is slightly different from the usual
stream classes, as a space is automatically inserted between arguments
of a \verb|<<|. This means that arguments can be easily given - eg
\begin{verbatim}
cmd << "plot" << 1.2 << 3.4 << "\n";
\end{verbatim}
will perform the command {\tt plot 1.2 3.4}. In order to build up a
symbol from multiple arguments, use the append operator. So
\begin{verbatim}
cmd << "plot"|3 << "\n";
\end{verbatim} 
executes the command {\tt plot3}.

The full class is given by:
\begin{verbatim}
class tclcmd
{
 public:
  char *result; /* The result of the previous command is placed here */
  tclcmd();
  ~tclcmd();
  tclcmd& operator<<(char* cmd);
  tclcmd& operator<<(double x);
  tclcmd& operator<<(int x); 
  tclcmd& operator<<(unsigned x);
  tclcmd& operator<<(tclvar x);
  tclcmd& operator|(char* cmd);
  tclcmd& operator|(double x);
  tclcmd& operator|(int x); 
  tclcmd& operator|(unsigned x);
  tclcmd& operator|(tclvar x);
};
\end{verbatim}

\subsection{tclreturn}

The {\tt tclreturn} class is used to supply a return value to a TCL
command. It inherits from {\tt ostrstream}, so has exactly the same
behaviour as that class. When this variable goes out of scope, the
contents of the stream buffer is written as the TCL return value. It is
an error to have more than one {\tt tclreturn} in a TCL command, the
result is undefined in this case.

\subsection{tclindex}

This class is used to index through a TCL array. TCL arrays can have
arbitrary strings for indices, and are not necessarily ordered. The
class has three main methods: {\tt start}, which takes a {\tt tclvar}
variable that refers to the array, and returns the first element in
the array; {\tt incr}, which returns the next element of the array;
and {\tt last} which returns 1 if the last element has been read. A
sample usage might be (for computing the product of all elements in
the array {\tt dims}):
\begin{verbatim}
tclvar dims="dims";
tclindex idx;
for (ncells *= (int)idx.start(dims); !idx.last(); ncells *= (int)idx.incr() );
\end{verbatim}

The class definition is given by:
\begin{verbatim}
class tclindex  
{
public:
  tclindex();
  tclindex(tclindex&);
  ~tclindex() {done();}
  tclvar start(tclvar&);
  inline tclvar incr();
  tclvar incr(tclvar& x) {incr();}  
  int last();
};
\end{verbatim}





